\documentclass[10pt,a4paper]{amsart}
\usepackage{amsmath,amssymb,amsthm,graphicx}
% Retained preamble for the Stacks Project chapter morphisms.tex (Morphisms of Schemes).
% The chapter file pinned-original-source.tex carries only \input{preamble} on its first line, so no
% preamble text lives inside the chapter itself.  The Stacks Project uses one single shared preamble for
% every chapter; the snapshot stages it verbatim in sources/stacks-chow/preamble.tex.  The block below is
% copied mechanically, line for line, from that file's lines 14-18 (the Xy-pic setup) and lines 154-217
% (the native theorem declarations together with the mathematical macros).  Nothing was added, removed or
% redefined.  The wrapper class and the amsmath/amssymb/amsthm/graphicx base packages are supplied by the
% assembler, so the class/\externaldocument part of the shared preamble is not repeated here.
% For commutative diagrams we use Xy-pic
\usepackage[all]{xy}

% We use 2cell for 2-commutative diagrams.
\xyoption{2cell}
% Theorem environments.
%
\theoremstyle{plain}
\newtheorem{theorem}[subsection]{Theorem}
\newtheorem{proposition}[subsection]{Proposition}
\newtheorem{lemma}[subsection]{Lemma}

\theoremstyle{definition}
\newtheorem{definition}[subsection]{Definition}
\newtheorem{example}[subsection]{Example}
\newtheorem{exercise}[subsection]{Exercise}
\newtheorem{situation}[subsection]{Situation}

\theoremstyle{remark}
\newtheorem{remark}[subsection]{Remark}
\newtheorem{remarks}[subsection]{Remarks}

\numberwithin{equation}{subsection}

% Macros
%
\def\lim{\mathop{\mathrm{lim}}\nolimits}
\def\colim{\mathop{\mathrm{colim}}\nolimits}
\def\Spec{\mathop{\mathrm{Spec}}}
\def\Hom{\mathop{\mathrm{Hom}}\nolimits}
\def\Ext{\mathop{\mathrm{Ext}}\nolimits}
\def\SheafHom{\mathop{\mathcal{H}\!\mathit{om}}\nolimits}
\def\SheafExt{\mathop{\mathcal{E}\!\mathit{xt}}\nolimits}
\def\Sch{\mathit{Sch}}
\def\Mor{\mathop{\mathrm{Mor}}\nolimits}
\def\Ob{\mathop{\mathrm{Ob}}\nolimits}
\def\Sh{\mathop{\mathit{Sh}}\nolimits}
\def\NL{\mathop{N\!L}\nolimits}
\def\CH{\mathop{\mathrm{CH}}\nolimits}
\def\proetale{{pro\text{-}\acute{e}tale}}
\def\etale{{\acute{e}tale}}
\def\QCoh{\mathit{QCoh}}
\def\Ker{\mathop{\mathrm{Ker}}}
\def\Im{\mathop{\mathrm{Im}}}
\def\Coker{\mathop{\mathrm{Coker}}}
\def\Coim{\mathop{\mathrm{Coim}}}

% Boxtimes
%
\DeclareMathSymbol{\boxtimes}{\mathbin}{AMSa}{"02}

%
% Macros for moduli stacks/spaces
%
\def\QCohstack{\mathcal{QC}\!\mathit{oh}}
\def\Cohstack{\mathcal{C}\!\mathit{oh}}
\def\Spacesstack{\mathcal{S}\!\mathit{paces}}
\def\Quotfunctor{\mathrm{Quot}}
\def\Hilbfunctor{\mathrm{Hilb}}
\def\Curvesstack{\mathcal{C}\!\mathit{urves}}
\def\Polarizedstack{\mathcal{P}\!\mathit{olarized}}
\def\Complexesstack{\mathcal{C}\!\mathit{omplexes}}
% \Pic is the operator that assigns to X its picard group, usage \Pic(X)
% \Picardstack_{X/B} denotes the Picard stack of X over B
% \Picardfunctor_{X/B} denotes the Picard functor of X over B
\def\Pic{\mathop{\mathrm{Pic}}\nolimits}
\def\Picardstack{\mathcal{P}\!\mathit{ic}}
\def\Picardfunctor{\mathrm{Pic}}
\def\Deformationcategory{\mathcal{D}\!\mathit{ef}}

% --- appended 2026-10-12 (Chief): the project own definitions of the reference/slogan/history
% display environments, copied mechanically from the shared Stacks preamble lines 8-12
% (\usepackage{verbatim} plus the three \newenvironment lines, which the chapter preamble keeps
% verbatim from the project preamble). Nothing added, removed or redefined: these environments
% carry no mathematical content and the project itself typesets them as comments. Without these
% three lines a counted statement whose own text uses them fails with Environment slogan undefined.
\usepackage{verbatim}
\newenvironment{reference}{\comment}{\endcomment}
\newenvironment{slogan}{\comment}{\endcomment}
\newenvironment{history}{\comment}{\endcomment}
\usepackage{cleveref}
\crefname{originallemma}{original lemma}{original lemmas}
\Crefname{originallemma}{Original lemma}{Original lemmas}
\makeatletter
\newlabel{algebra-lemma-Noetherian-irreducible-components}{{00FR}{source}{}{sourcepointer-1}{}}
\newlabel{algebra-lemma-Noetherian-irreducible-components@cref}{{[originallemma][0][]00FR}{[1][0][]source}{}{}{}}
\makeatother
\crefname{originallemma}{original lemma}{original lemmas}
\Crefname{originallemma}{Original lemma}{Original lemmas}
\makeatletter
\newlabel{algebra-lemma-characterize-reduced-ring-normal}{{030C}{source}{}{sourcepointer-2}{}}
\newlabel{algebra-lemma-characterize-reduced-ring-normal@cref}{{[originallemma][0][]030C}{[1][0][]source}{}{}{}}
\makeatother
\crefname{originallemma}{original lemma}{original lemmas}
\Crefname{originallemma}{Original lemma}{Original lemmas}
\makeatletter
\newlabel{algebra-lemma-integral-closure-localize}{{0307}{source}{}{sourcepointer-3}{}}
\newlabel{algebra-lemma-integral-closure-localize@cref}{{[originallemma][0][]0307}{[1][0][]source}{}{}{}}
\makeatother
\crefname{originallemma}{original lemma}{original lemmas}
\Crefname{originallemma}{Original lemma}{Original lemmas}
\makeatletter
\newlabel{algebra-lemma-silly}{{00DS}{source}{}{sourcepointer-4}{}}
\newlabel{algebra-lemma-silly@cref}{{[originallemma][0][]00DS}{[1][0][]source}{}{}{}}
\makeatother
\crefname{originallemma}{original lemma}{original lemmas}
\Crefname{originallemma}{Original lemma}{Original lemmas}
\makeatletter
\newlabel{lemma-nagata-normalization-finite}{{03GR}{source}{}{sourcepointer-5}{}}
\newlabel{lemma-nagata-normalization-finite@cref}{{[originallemma][0][]03GR}{[1][0][]source}{}{}{}}
\makeatother
\crefname{originallemma}{original lemma}{original lemmas}
\Crefname{originallemma}{Original lemma}{Original lemmas}
\makeatletter
\newlabel{lemma-normalization-is-normalization}{{0BXA}{source}{}{sourcepointer-6}{}}
\newlabel{lemma-normalization-is-normalization@cref}{{[originallemma][0][]0BXA}{[1][0][]source}{}{}{}}
\makeatother
\crefname{originallemma}{original lemma}{original lemmas}
\Crefname{originallemma}{Original lemma}{Original lemmas}
\makeatletter
\newlabel{lemma-reach-points-scheme-theoretic-image}{{02JQ}{source}{}{sourcepointer-7}{}}
\newlabel{lemma-reach-points-scheme-theoretic-image@cref}{{[originallemma][0][]02JQ}{[1][0][]source}{}{}{}}
\makeatother
\title{A one-dimensional point of a normalization outside the image of the source has discrete valuation ring}
\author{The Stacks Project Authors}
\date{stacks-project:morphisms.tex:local-snapshot-20261012; GNU FDL 1.2 or later}
\begin{document}
\maketitle
\noindent\textit{Editorial note.} All delivered mathematical text is the author's own, copied from the pinned Stacks Project chapter file pinned-original-source.tex; no statement or proof step was rewritten, reordered or removed.\par\medskip
\noindent\textit{Editorial note.} Following the multi-claim convention for this source, source\_url is the source-level URL https://stacks.math.columbia.edu/ and the source label lemma-relative-normalization-normal-codim-1 is carried in claim\_key and here rather than in a per-claim tag URL.\par\medskip
\noindent\textit{Editorial note.} The preamble is the chapter's shared Stacks Project preamble, retained as sources/stacks-morphisms/retained-preamble.tex: a line-for-line excerpt of the shared Stacks preamble staged with this snapshot. No macro was invented or redefined, and this card adds no preamble text of its own.\par\medskip
\noindent\textit{Editorial note.} Statement span 14292-14302 and proof span 14304-14348 come from the halfA pool card operations/pool-cards/stacks-morphisms-halfA.json and were re-read in the pinned chapter file before this card was written.\par\medskip
\noindent\textit{Editorial note.} Re-checked by hand before carding: this candidate was caught by the lane's thin-proof prefilter, and reading the original shows the hit comes only from the closing phrase 'Some details omitted.' while the proof is a full reduction argument (finiteness of the normalization, prime avoidance, vanishing-locus analysis, integral-closure computation). Per the fleet rule of 2026-10-14 the prefilter hit on a proof of more than 800 bytes was re-read by hand, and the reason is recorded here. The author's closing phrase is kept verbatim in the counted span.\par\medskip
\noindent\textit{Editorial note.} 7 cross-references occur inside the counted spans; they are listed in external\_references as navigation pointers only.\par\medskip
\noindent\textit{Allowed standard prerequisites (editorial).}\begin{itemize}
\item For a Nagata base and a finite type morphism with reduced source the normalization in the source is finite, cited from the same chapter. Cited at proof lines 14299-14301. Makes the normalization affine and finite over an affine base, so that the situation becomes a finite algebra over a Noetherian ring.
\item The idempotence of the normalization, established earlier in the same chapter (and by a previous card of this lane), which lets the base be replaced by the normalization itself. Cited at proof lines 14301-14303. Reduces to the case where base and normalization coincide.
\item Prime avoidance together with the finiteness of the set of minimal primes of a Noetherian ring, and the description of the closure of the image of a closed subset as the specializations of its points. Cited at proof lines 14306-14309 and 14318-14320. Produce the element avoiding the minimal primes and show that the closure of the image of its vanishing locus avoids the point.
\item The compatibility of integral closure with localization, and the characterization of a reduced ring as normal via its local behaviour, cited from the algebra chapter. Cited at proof lines 14333-14335 and 14342-14345. Turn the equality between the ring and the integral closure in a localization into normality of the local ring at the point.
\end{itemize}\medskip
\ifdefined\hypertarget\hypertarget{sourcepointer-1}{}\fi
\noindent\textit{Source reference.} Original lemma 00FR: focused-seeds/AGP/sources/stacks-snapshot-20261011/tags.txt, source lines 581--581. Same-chapter reference used in the retained author text at source line(s) 13632, 14321, 15637. It is supplied only as a navigation pointer so that the retained text compiles; the counted item delivers the author's own statement and the author's complete proof, and no part of the referenced item is copied in or substituted for the counted argument.\par\medskip
\ifdefined\hypertarget\hypertarget{sourcepointer-2}{}\fi
\noindent\textit{Source reference.} Original lemma 030C: focused-seeds/AGP/sources/stacks-snapshot-20261011/tags.txt, source lines 3943--3943. Same-chapter reference used in the retained author text at source line(s) 14346, 14454. It is supplied only as a navigation pointer so that the retained text compiles; the counted item delivers the author's own statement and the author's complete proof, and no part of the referenced item is copied in or substituted for the counted argument.\par\medskip
\ifdefined\hypertarget\hypertarget{sourcepointer-3}{}\fi
\noindent\textit{Source reference.} Original lemma 0307: focused-seeds/AGP/sources/stacks-snapshot-20261011/tags.txt, source lines 3938--3938. Same-chapter reference used in the retained author text at source line(s) 13874, 14339, 14493. It is supplied only as a navigation pointer so that the retained text compiles; the counted item delivers the author's own statement and the author's complete proof, and no part of the referenced item is copied in or substituted for the counted argument.\par\medskip
\ifdefined\hypertarget\hypertarget{sourcepointer-4}{}\fi
\noindent\textit{Source reference.} Original lemma 00DS: focused-seeds/AGP/sources/stacks-snapshot-20261011/tags.txt, source lines 508--508. Same-chapter reference used in the retained author text at source line(s) 14320. It is supplied only as a navigation pointer so that the retained text compiles; the counted item delivers the author's own statement and the author's complete proof, and no part of the referenced item is copied in or substituted for the counted argument.\par\medskip
\ifdefined\hypertarget\hypertarget{sourcepointer-5}{}\fi
\noindent\textit{Source reference.} Original lemma 03GR: focused-seeds/AGP/sources/stacks-snapshot-20261011/tags.txt, source lines 4538--4538. Same-chapter reference used in the retained author text at source line(s) 14307. It is supplied only as a navigation pointer so that the retained text compiles; the counted item delivers the author's own statement and the author's complete proof, and no part of the referenced item is copied in or substituted for the counted argument.\par\medskip
\ifdefined\hypertarget\hypertarget{sourcepointer-6}{}\fi
\noindent\textit{Source reference.} Original lemma 0BXA: focused-seeds/AGP/sources/stacks-snapshot-20261011/tags.txt, source lines 14973--14973. Same-chapter reference used in the retained author text at source line(s) 14310. It is supplied only as a navigation pointer so that the retained text compiles; the counted item delivers the author's own statement and the author's complete proof, and no part of the referenced item is copied in or substituted for the counted argument.\par\medskip
\ifdefined\hypertarget\hypertarget{sourcepointer-7}{}\fi
\noindent\textit{Source reference.} Original lemma 02JQ: focused-seeds/AGP/sources/stacks-snapshot-20261011/tags.txt, source lines 3345--3345. Same-chapter reference used in the retained author text at source line(s) 1367, 4438, 10265, 14326. It is supplied only as a navigation pointer so that the retained text compiles; the counted item delivers the author's own statement and the author's complete proof, and no part of the referenced item is copied in or substituted for the counted argument.\par\medskip
\begin{lemma}
\label{lemma-relative-normalization-normal-codim-1}
Let $f : Y \to X$ be a finite type morphism of schemes with $Y$ reduced
and $X$ Nagata. Let $X'$ be the normalization of $X$ in $Y$.
Let $x' \in X'$ be a point such that
\begin{enumerate}
\item $\dim(\mathcal{O}_{X', x'}) = 1$, and
\item the fibre of $Y \to X'$ over $x'$ is empty.
\end{enumerate}
Then $\mathcal{O}_{X', x'}$ is a discrete valuation ring.
\end{lemma}

\begin{proof}
We can replace $X$ by an affine neighbourhood of the image of $x'$.
Hence we may assume $X = \Spec(A)$ with $A$ Nagata.
By Lemma \ref{lemma-nagata-normalization-finite}
the morphism $X' \to X$ is finite. Hence we can write
$X' = \Spec(A')$ for a finite $A$-algebra $A'$.
By Lemma \ref{lemma-normalization-is-normalization}
after replacing $X$ by $X'$
we reduce to the case described in the next paragraph.

\medskip\noindent
The case $X = X' = \Spec(A)$ with $A$ Noetherian.
Let $\mathfrak p \subset A$ be the prime ideal
corresponding to our point $x'$. Choose $g \in \mathfrak p$
not contained in any minimal prime of $A$ (use prime avoidance
and the fact that $A$ has finitely many minimal primes, see
Algebra, Lemmas \ref{algebra-lemma-silly} and
\ref{algebra-lemma-Noetherian-irreducible-components}).
Set $Z = f^{-1}V(g) \subset Y$; it is a closed subscheme of $Y$.
Then $f(Z)$ does not contain any generic point by choice of $g$
and does not contain $x'$ because $x'$ is not in the image of $f$.
The closure of $f(Z)$ is the set of specializations of points
of $f(Z)$ by Lemma \ref{lemma-reach-points-scheme-theoretic-image}.
Thus the closure of $f(Z)$ does not contain $x'$ because
the condition $\dim(\mathcal{O}_{X', x'}) = 1$ implies only
the generic points of $X = X'$ specialize to $x'$.
In other words, after replacing $X$ by an affine open
neighbourhood of $x'$ we may assume that $f^{-1}V(g) = \emptyset$.
Thus $g$ maps to an invertible global function on $Y$ and
we obtain a factorization
$$
A \to A_g \to \Gamma(Y, \mathcal{O}_Y)
$$
Since $X = X'$ this implies that $A$ is equal to the integral
closure of $A$ in $A_g$. By
Algebra, Lemma \ref{algebra-lemma-integral-closure-localize}
we conclude that $A_\mathfrak p$ is the integral closure
of $A_\mathfrak p$ in $A_\mathfrak p[1/g]$.
By our choice of $g$, since $\dim(A_\mathfrak p) = 1$
and since $A$ is reduced we see that $A_\mathfrak p[1/g]$
is a finite product of fields (the product of the residue fields
of the minimal primes contained in $\mathfrak p$). Hence $A_\mathfrak p$ is
normal (Algebra, Lemma \ref{algebra-lemma-characterize-reduced-ring-normal})
and the proof is complete. Some details omitted.
\end{proof}

\end{document}
